% I. ПРЕДМЕТНА ОБЛАСТ И ЛИТЕРАТУРЕН ПРЕГЛЕД
% Покрива от учебната програма: характеристики на UML, понятиен модел, градивни
% елементи, връзки, диаграми, правила и общи механизми, моделиране на системи,
% жизнен цикъл на програмния продукт и UML.
\section{Предметна област и литературен преглед}
\label{sec:literature}

UML е графичен език за визуализиране, специфициране, конструиране и документиране на
артефактите на една програмна система. Той не е метод и не предписва процес на
разработка, а речник и правила за съчетаване на думите от този речник. Тази разлика е
съществена за настоящия проект: щом езикът е формален, а не илюстративен, той има
недвусмислена абстрактна структура и следователно подлежи на машинна обработка.
Официалната спецификация на езика~\cite{omg_uml251} задава именно тази структура под
формата на метамодел, а ръководствата на авторите на
езика~\cite{booch2005userguide,rumbaugh2004reference} го излагат в приложим за
проектант вид.

Понятийният модел на UML стъпва на три групи съставки. Първата са \term{градивните
елементи}: структурни (клас, интерфейс, компонент, възел), поведенчески
(взаимодействие, състояние) и групиращи (пакет), заедно с пояснителния елемент
(бележка). Втората са \term{връзките}: зависимост, асоциация, обобщение и реализация,
като асоциацията има отделни разновидности за агрегация и композиция, различаващи се по
това дали частта може да съществува без цялото. Третата са \term{правилата и общите
механизми}: спецификации, украшения, общи разделения и \term{механизмите за
разширяване}, тоест ограничения, етикетирани стойности и стереотипи. Настоящият проект
се интересува непосредствено от последните, защото именно чрез стереотип и етикетирана
стойност се обозначава кой клас участва в релационната схема и как.

Диаграмата на класове е основната структурна диаграма и е носител на речника на
системата. Класификаторът в нея се описва с име, атрибути, операции и отговорности,
като видимостта и обхватът се отбелязват върху всеки член. Практиката за откриване на
този речник е по-стара от самия UML: картите клас, отговорност и сътрудничество, познати
като CRC карти, са предложени от Beck и Cunningham~\cite{beck1989crc} точно като начин
класовете да бъдат намерени преди да бъдат нарисувани. За компактно въведение в
приложната част на езика е използван и~\cite{fowler2003distilled}.

Отношението между обектния модел и релационния модел на данните е самостоятелна тема със
собствена литература. Релационната схема се описва чрез таблици, ключове и връзки между
тях и има свой понятиен предшественик в модела същност и връзка на
Chen~\cite{chen1976er}. Съответствието между клас и таблица не е взаимно еднозначно:
наследяването няма пряк аналог в релационния модел, а асоциацията се изразява чрез външен
ключ или чрез свързваща таблица според кратността. Систематичен преглед на утвърдените
стратегии за това съответствие дава~\cite{ambler2003agiledb}. Тези стратегии са
непосредствената основа на правото инженерство в настоящия проект.

Мястото на двупосочното инженерство в жизнения цикъл на програмния продукт следва
таксономията на Chikofsky и Cross~\cite{chikofsky1990reverse}. Обратното инженерство
поддържа документацията вярна, без да изисква дисциплина от разработчика, и е приложимо
и към чужд код. Правото инженерство пренася решение, взето на нивото на модела, надолу
до реализацията, без ръчно преписване. Стойността на двете заедно е в третото им
следствие: щом и двете посоки са автоматизирани, разминаването между модел и код става
измеримо, а не просто подозирано. \TODO{при разширяване на раздела да се добави
преглед на инструментите за проверка на съответствие между модел и код}
